\section*{第 \ref{ch:b1:titration-methods} 章——滴定方法与滴定曲线}

\begin{solution}{exo:b1:titration-methods:1}
\ce{H3O+ + OH- -> 2H2O}, $K = 1/K_e = 10^{14.00}$. $C = 12.4 \times
0.100/20.0 = \qty{0.0620}{mol/L}$.
\end{solution}

\begin{solution}{exo:b1:titration-methods:2}
化学计量点时，溶液为氯化铵，浓度为
\qty{0.050}{mol/L}：$\mathrm{pH} = \frac12(9.25 - \log 0.050) = 5.28$。
甲基红（变色区间 3.8--5.8）包含 5.28 这一 pH 值；甲基橙的变色区间为 2.5--4.5，只有在化学计量点之后随 pH 下降才能达到，会变色过晚。
\end{solution}

\begin{solution}{exo:b1:titration-methods:3}
\ce{MnO4- + 5Fe^2+ + 8H+ -> Mn^2+ + 5Fe^3+ + 4H2O};
$C = 5 \times 0.0200 \times 12.6/10.0 = \qty{0.126}{mol/L}$.
\end{solution}

\begin{solution}{exo:b1:titration-methods:4}
钙与镁的总浓度 $C = 0.0100 \times 14.2/50.0 = \qty{2.84e-3}{mol/L}$，
两者均在滴定终点前反应，终点只给出二者之和：
即同时滴定。
\end{solution}

\begin{solution}{exo:b1:titration-methods:5}
0~mL：弱酸，$\mathrm{pH} = \frac12(4.76 + 1.00) = 2.88$。5.0~mL：
半等当点，4.76。10.0~mL：8.73
（\cref{prop:b1:titration-methods:ph-at-equivalence}）。15.0~mL：过量的
氢氧根 \qty{0.50}{mmol} 位于 \qty{25.0}{mL} 中，$[\ce{OH-}] = 0.020$，
$\mathrm{pH} = 12.30$。所有数据与计算曲线相符，偏差不超过 0.01。
\end{solution}

\begin{solution}{exo:b1:titration-methods:6}
$\mathrm{p}K_a$ 2.15 与 7.21 相差超过 4，7.21 与
12.34 亦然：有两个化学计量点，位于 10.0 和 \qty{20.0}{mL} 处。第一处为
\ce{H2PO4-} 溶液，$\mathrm{pH} \approx \frac12(2.15 + 7.21) = 4.68$；第二处为
\ce{HPO4^2-}，$\frac12(7.21 + 12.34) = 9.78$（精确曲线给出 4.71 和
9.66，公式忽略了稀释与水的作用）。第三步解离
（$\mathrm{p}K_a$ 为 12.34）与水的酸性过于接近，所加入的
氢氧根只能部分反应，pH 上升而没有突跃。
\end{solution}

\begin{solution}{exo:b1:titration-methods:7}
盐酸：$6.2 \times 0.100/10.0 = \qty{0.062}{mol/L}$；乙酸
：$(14.6 - 6.2) \times 0.100/10.0 = \qty{0.084}{mol/L}$。滴定至一半时，
乙酸已被滴定一半：$\mathrm{pH} = \mathrm{p}K_a = 4.76$。
\end{solution}

\begin{solution}{exo:b1:titration-methods:8}
$V_{\mathrm{eq}} = \qty{10.0}{mL}$。2.0~mL：$E = 0.77 + 0.059\log(2/8) =
0.73$~V；9.0~mL：$0.77 + 0.059\log 9 = 0.83$~V；11.0~mL：过量
高锰酸根 \qty{0.020}{mmol}，\ce{Mn^2+} \qty{0.200}{mmol}，$E = 1.51 +
\frac{0.059}{5}\log 0.10 = 1.50$~V；20.0~mL：1.51~V。
\end{solution}

\begin{solution}{exo:b1:titration-methods:9}
0~mL：$\kappa = (349.8 + 76.2) \times 0.0100 = \qty{4.26}{mS/cm}$。
10.0~mL：\ce{Na+} 与 \ce{Cl-} 浓度均为 $1.00/110 = \qty{9.09e-3}{mol/L}$，
$\kappa = (50.3 + 76.2) \times \num{9.09e-3} = 1.15$。20.0~mL：\ce{Na+}
为 $0.0167$，\ce{Cl-} 与 \ce{OH-} 分别为 $0.00833$~mol/L，$\kappa = 0.84 + 0.635 + 1.65
= 3.13$~mS/cm。斜率（忽略稀释，每 \qty{1}{mL} 加入
\qty{1.0e-3}{mol/L}）分别为 $(50.3 - 349.8) \times 10^{-3} = -0.30$ 与
$(50.3 + 198.3) \times 10^{-3} = +0.25$~mS/cm（每毫升）。
\end{solution}

\begin{solution}{exo:b1:titration-methods:10}
在 $V_{\mathrm{eq}} - v$ 处，过量的酸为 $C_Bv$；在 $V_{\mathrm{eq}} + v$
处，过量的碱为 $C_Bv$。同一体积中，前者的 $[\ce{H3O+}]$ 与后者的
$[\ce{OH-}]$ 相等，故 $\mathrm{pH}(V_{\mathrm{eq}} - v) +
\mathrm{pH}(V_{\mathrm{eq}} + v) = 14$，即关于 (V$_{\mathrm{eq}}$,
7) 对称。突跃：在 \qty{20}{mL} 中过量 \qty{1.0e-5}{mol}，pH 由 3.30 变至 10.70，变化 7.4
个单位。若除以 100：$\num{1.0e-7}$~mol 位于 \qty{20}{mL} 中，pH 由 5.3 变至 8.7，
变化 3.4 个单位：终点远不够敏锐。
\end{solution}

\begin{solution}{exo:b1:titration-methods:11}
化学计量点前：生成的 \ce{Fe^3+} 为 $= 5C_{\mathrm{Mn}}V$，剩余的 \ce{Fe^2+} 为 $=
5C_{\mathrm{Mn}}(V_{\mathrm{eq}} - V)$，故 $E = 0.77 + 0.059\log\frac{V}{V_{\mathrm{eq}}
- V}$。当 $V_{\mathrm{eq}}/2$ 时，$E = 0.77$。化学计量点后，过量的 \ce{MnO4-} $\propto
V - V_{\mathrm{eq}}$，\ce{Mn^2+} $\propto V_{\mathrm{eq}}$，$E = 1.51 +
\frac{0.059}{5}\log\frac{V - V_{\mathrm{eq}}}{V_{\mathrm{eq}}}$（pH 0）；当
$2V_{\mathrm{eq}}$ 时为 1.51~V。化学计量点处 $(0.77 + 5 \times 1.51)/6 =
1.39$~V。在 pH 1 时，高锰酸根电对的电势降低 $0.059 \times 8/5 =
0.094$~V：$E_{\mathrm{eq}} = (0.77 + 5 \times 1.416)/6 = 1.31$~V。
\end{solution}

\begin{solution}{exo:b1:titration-methods:12}
\ce{NH4+ + OH- -> NH3 + H2O} 的 $K = 10^{14.00 - 9.25} = 10^{4.75}$，故
化学计量点附近反应不完全，突跃较小，
约在 pH 11 处（化学计量点 pH $14 - \frac12(4.75 + 1.30) = 11.0$）。返滴定：
加入过量的 $n_1$ 氢氧化钠，煮沸以逸出生成的氨，
再用盐酸滴定剩余的氢氧根
（$n_2$，强酸与强碱的突跃明显）：$n(\ce{NH4+}) = n_1 - n_2$。
\end{solution}

\begin{solution}{pb:b1:titration-methods:1}
\textbf{1.} \ce{C6H6O6 + 2H+ + 2e- -> C6H8O6}.
\textbf{2.} \ce{C6H8O6 + I2 -> C6H6O6 + 2I- + 2H+}.
\textbf{3.} 0 和 $-$I。
\textbf{4.} \ce{I2 + 2S2O3^2- -> 2I- + S4O6^2-}, $\log K = 2(0.53 - 0.02)/0.059
= 17$.
\textbf{5.} 蓝色需要碘存在；蓝色随最后一点
碘消失，即在化学计量点。
\textbf{6.} 空气中的损失会使结果偏低；一次加入并过量时，
碘能迅速而完全地氧化抗坏血酸。
\textbf{7.} 整个滴定过程中，反应快速、定量，
滴定剂缓慢加入时无空气氧化，终点位于
碘首次过量处。
\textbf{8.} 过量碘 $n_R$；与抗坏血酸反应；
用硫代硫酸盐滴定剩余的碘。
\textbf{9.} 否则会有部分抗坏血酸残留，而碘
的过量为零，便无法说明残留多少。此处
有 \qty{2.775e-4}{mol} 发生反应（总量 \qty{5.000e-4}{mol}），确实过量。
\textbf{10.} 硫代硫酸盐溶液随时间缓慢变化；其
浓度直接影响结果。
\textbf{11.} 10.
\textbf{12.} 一支 \qty{20.00}{mL} 容量移液管。
\textbf{13.} $20.00 \times 10^{-3} \times 0.0250 = \qty{5.000e-4}{mol}$.
\textbf{14.} $8.90 \times 10^{-3} \times 0.0500 = \qty{4.450e-4}{mol}$.
\textbf{15.} 一半：\qty{2.225e-4}{mol}。
\textbf{16.} $5.000 - 2.225 = \qty{2.775e-4}{mol}$.
\textbf{17.} 一比一：\qty{2.775e-4}{mol}。
\textbf{18.} $\times 10$: \qty{2.775e-3}{mol}, $\times 176.0$:
\qty{0.488}{g}.
\textbf{19.} 比标示量低约 \qty{2}{\percent}。
\textbf{20.} $n = \bigl(C_IV_I - \frac12C_TV_T\bigr)\,V_{\mathrm{flask}}/V_{\mathrm{aliquot}}$.
\textbf{21.} 引入的碘：移液管 $0.03/20.00 = \qty{0.15}{\percent}$
与浓度 \qty{0.2}{\percent}，合计约 \qty{0.25}{\percent}，即
\qty{1.3e-6}{mol}。过量部分：滴定管 $0.05/8.90 = \qty{0.56}{\percent}$ 与
浓度 \qty{0.2}{\percent}，合计约 \qty{0.60}{\percent}，即
\qty{1.3e-6}{mol}.
\textbf{22.} 绝对不确定度合成，而两者之差
小于任一量：$\sqrt{1.3^2 + 1.3^2} \times 10^{-6} =
\qty{1.8e-6}{mol}$，占 \qty{2.775e-4}{mol} 的 \qty{0.66}{\percent}，超过
任一项。
\textbf{23.} 计入容量瓶（\qty{0.1}{\percent}）与分取移液管
(\qty{0.2}{\percent}): $\sqrt{0.66^2 + 0.1^2 + 0.2^2} = \qty{0.70}{\percent}$.
\textbf{24.} $0.488 \pm 0.003$~g：\textbf{每片含抗坏血酸
\qty{0.49}{g}}。
\end{solution}
